\documentclass[11pt]{article}
\usepackage{amsmath}
\usepackage[hidelinks]{hyperref}

\title{Ropes for Text Editors:\\ A Short Example Paper}
\author{An Example Author}
\date{}

\begin{document}
\maketitle

\begin{abstract}
This short paper exists to demonstrate a review workflow for LaTeX documents.
It describes ropes, a tree-shaped representation of long strings that text
editors use to make insertions and deletions cheap. Its content is
deliberately ordinary.
\end{abstract}

\section{Introduction}
\label{sec:introduction}
A text editor keeps the document it is editing in memory, and every key
press changes it. The simplest representation, a single contiguous array of
characters, makes reading fast but makes editing slow: inserting one
character near the start of a large file moves every character after it.

Editors have answered this problem in several ways. A gap buffer keeps an
empty region at the cursor, so that typing at the cursor is cheap while
jumping elsewhere is not~\cite{gap}. A piece table records edits as a list
of references into the original file and an append-only buffer. A rope,
the subject of this paper, stores the text in the leaves of a balanced
binary tree~\cite{boehm}.

Section~\ref{sec:design} describes how a rope represents text and how the
basic operations work. Section~\ref{sec:evaluation} compares ropes with the
alternatives on a few common editing patterns.

\section{Design}
\label{sec:design}
A rope is a binary tree whose leaves hold short strings. Reading the leaves
from left to right yields the whole document. Each node $v$ carries a
weight $w(v)$, the number of characters in its left subtree:
\begin{equation}
  w(v) = |\mathrm{left}(v)|.
  \label{eq:weight}
\end{equation}
Weights are the only extra information a rope needs. Every operation below
reads them on the way down the tree and updates them on the way back up,
so keeping them correct costs nothing beyond the nodes an operation already
touches.

\subsection{Finding a position}
\label{sec:finding}
To find the character at position $i$, we start at the root and compare $i$
with the weight of the current node. If $i < w(v)$, we continue in the left
subtree; otherwise we continue in the right subtree with $i - w(v)$. When we
reach a leaf, $i$ is an offset into its string. A rope of depth $d$ answers
this query in $O(d)$ steps.

We call a rope balanced when its depth $d$ and its length $n$ satisfy
\begin{equation}
  n \ge F_{d+2},
  \label{eq:balance}
\end{equation}
where $F_k$ is the $k$-th Fibonacci number~\cite{boehm}. Since $F_k$ grows
like $\varphi^k$ with $\varphi = (1+\sqrt{5})/2$, condition~\eqref{eq:balance}
bounds the depth of a balanced rope by $\log_\varphi n + O(1)$, so finding a
position takes logarithmic time.

\subsection{Editing}
\label{sec:editing}
Splitting walks down the same path as finding a position and creates new
nodes only along that path. Concatenation creates a single new root, whose
weight is the length of its left argument.
In order to be able to insert a string at position $i$, we first split the rope at $i$ into two ropes, and then we concatenate the left part, the new string, and the right part.
Neither operation copies the text stored in the leaves, which is why ropes
stay fast for large documents.

\begin{figure}[t]
  \centering
  \fbox{split at $i$} \enspace$\rightarrow$\enspace
  \fbox{left part} \enspace$+$\enspace \fbox{new text} \enspace$+$\enspace \fbox{right part}
  \enspace$\rightarrow$\enspace \fbox{concatenate}
  \caption{Inserting text into a rope.}
  \label{fig:insertion}
\end{figure}


Each edit therefore costs a split and at most two concatenations. Over a
sequence of $k$ edits, where the document has $n_j$ characters before the
$j$-th edit, the total cost is
\begin{align}
  T(k) &= \sum_{j=1}^{k} O(\log n_j) \label{eq:cost} \\
       &\le k \cdot O(\log n_k).
\end{align}
The bound assumes that the rope stays balanced between edits.

Deletion works in the same way: split at both ends of the deleted range and
concatenate the outer parts. Repeated edits can unbalance the tree, so an
editor occasionally rebuilds it, which restores condition~\eqref{eq:balance}
in time linear in the number of leaves.
Rebalancing is cheap enough to run after every few hundred edits.

The size of the leaves is a tuning parameter. Very short leaves waste memory
on internal nodes, while very long leaves make every split copy more text
than necessary. Leaves of a few hundred bytes work well in practice, because
they fit in a handful of cache lines and keep the tree shallow even for
documents of many megabytes.

Because old nodes are never modified, a rope also gives the editor an undo
history almost for free: each version of the document is simply a pointer
to an older root.

\section{Evaluation}
\label{sec:evaluation}
We compare the three representations on three patterns: typing at the end
of a file, typing in the middle of a large file, and replacing every
occurrence of a word. Table~\ref{tab:costs} summarizes the cost of a single
edit at distance $\delta$ from the previous one in a file of $n$ characters.

\begin{table}[h]
  \centering
  \begin{tabular}{lccc}
    & typing at the end & typing in the middle & reading \\
    \hline
    gap buffer  & $O(1)$       & $O(\delta)$ & $O(1)$ \\
    piece table & $O(1)$       & $O(p)$      & $O(p)$ \\
    rope        & $O(\log n)$ & $O(\log n)$ & $O(\log n)$ \\
  \end{tabular}
  \caption{Cost of one edit; $p$ is the number of pieces.}
  \label{tab:costs}
\end{table}

When typing at the end of a file, all three representations perform well.
The gap buffer is fastest, because the gap is already at the cursor and
each key press writes one byte.

When typing in the middle of a large file, the gap buffer must first move
the gap, which costs time proportional to the distance from the previous
cursor position. The piece table and the rope are unaffected by the
distance.

Replacing every occurrence of a word makes many edits spread across the
document. The rope handles each edit in logarithmic time and keeps every
intermediate version, while the piece table accumulates many small pieces
that slow down later reads until it is compacted.

Taken together, these patterns suggest that ropes are a good default for
editors that open large files, and that gap buffers remain attractive for
small files and simple editors.


\section{Conclusion}
\label{sec:conclusion}
Ropes trade a little constant overhead for edits whose cost does not grow
with the length of the document. For an editor, which spends most of its
time making small edits to large buffers, this is usually the right trade.

\begin{thebibliography}{9}
\bibitem{boehm} H.-J. Boehm, R. Atkinson, and M. Plass.
  Ropes: an alternative to strings.
  \emph{Software: Practice and Experience}, 25(12), 1995.
\bibitem{gap} C. Finseth. \emph{The Craft of Text Editing}. Springer, 1991.
\end{thebibliography}

\end{document}
